%% toc: P15 157 157.1,157.2,157.3,157.4
% Poisson and geometric expectation and moment calculations (Durrett, Examples 1.6.13 and 1.6.14).
\chapter{Poisson and Geometric Moment Calculations}
\label{sec:p15-w02b-p15-001-poisson-geometric}

This chapter computes the expectations and moments of the Poisson and
geometric laws, following Examples~1.6.13 and~1.6.14 of Durrett's author
version of \emph{Probability: Theory and Examples}, dated 11 January 2019
\cite[pp.~34--35]{DurrettPTEAuthor2019}. The purpose is to show how the
integral of Chapter~\ref{sec:p15-integration-expectation} turns into a
series computation for a variable with countably many values, and to
record the resulting moment formulas with every interchange of limits
justified.

\emph{Inputs.} A real random variable whose law is Poisson with
parameter \(\lambda\), or geometric with success probability \(p\).
The integral, expectation and variance are those of
\cref{def:p15-nonnegative-integral,def:p15-expectation-integral}.
\emph{Conclusions.} Every descending factorial moment of each law in
closed form; the mean, the second moment and the variance; and, for the
geometric law, its realization as the index of the first occurrence among
independent events of equal probability.

Two features of the source route are changed here. Durrett obtains the
series behind the geometric moments by differentiating the geometric
series term by term and refers to an appendix result for the exchange
\cite[Example~A.5.6, pp.~472--473]{DurrettPTEAuthor2019}. The proofs
below use no differentiation: they regroup a nonnegative double series,
which \cref{thm:p15-countable-atomic} permits without any convergence
hypothesis. Second, the geometric factorial moments are stated for every
order \(r\); Durrett states only the orders one and two, and the higher
orders are an extension made here.

\section{Laws, descending products and the series tools}
\label{sec:p15-w02b-p15-001-laws}

\begin{hbdefinition}{def:p15-w02b-p15-001-laws}{Poisson and geometric laws and descending products}
Let \(\lambda>0\). A real random variable \(X\) has the \emph{Poisson
distribution with parameter \(\lambda\)} if
\[
 \mathbb P(X=k)=e^{-\lambda}\frac{\lambda^k}{k!}
 \qquad(k=0,1,2,\ldots).
\]
These masses sum to one by \cref{thm:p15-stirling-poisson-premises},
so \(\mathbb P(X\in\{0,1,2,\ldots\})=1\). The parameter-zero law, the
point mass at zero, is the convention of
\cref{def:p15-stirling-poisson-laws}.

Let \(0<p<1\) and write \(q\defeq1-p\). A real random variable \(N\)
has the \emph{geometric distribution with success probability \(p\)} if
\[
 \mathbb P(N=k)=p(1-p)^{k-1}=pq^{k-1}\qquad(k=1,2,\ldots).
\]
These masses sum to one by part~(1) of
\cref{thm:p15-w02b-p15-001-geometric-moments}, so
\(\mathbb P(N\in\{1,2,\ldots\})=1\). The endpoint \(p=1\) is the point
mass at one, as in \cref{def:p15-collector-and-cycle-models}.

For a real number \(x\) and an integer \(k\geq1\), the
\emph{descending product of order \(k\)} is
\(x(x-1)\cdots(x-k+1)\), the product of the \(k\) factors \(x-i\),
\(0\leq i\leq k-1\). Its expectation, when defined, is the
\emph{factorial moment of order \(k\)}. For an integer \(j\geq0\) the
descending product is zero when \(j<k\), because the factor with
\(i=j\) vanishes, and it equals \(j!/(j-k)!\) when \(j\geq k\). On the
nonnegative integers it is therefore nonnegative.
For integers \(0\leq r\leq n\),
\(\binom nr\defeq n!/(r!\,(n-r)!)\).
\index{Poisson distribution!moments}\index{geometric distribution}
\index{factorial moment}
\end{hbdefinition}

\begin{hblemma}{thm:p15-w02b-p15-001-series-tools}{proved=annex:p15-w02b-p15-001-series-tools}{Expectations as series and the negative-binomial series}
\par\nobreak
\begin{enumerate}
\item Let \(X\) be a real random variable and \(S\subseteq\mathbb R\)
      a countable set with \(\mathbb P(X\in S)=1\). For every Borel
      \(h:\mathbb R\to[0,\infty)\),
      \[
       \mathbb E[h(X)]=\sum_{s\in S}h(s)\,\mathbb P(X=s)
      \]
      in \([0,\infty]\). If instead \(h:\mathbb R\to\mathbb R\) is
      Borel and \(h\geq0\) on \(S\), the random
      variable \(h(X)\) is nonnegative a.s., its expectation is defined,
      and the same identity holds.
\item Let \(0\leq q<1\) and \(p=1-q\). For every integer \(r\geq0\),
      \[
       \sum_{m\geq0}\binom{m+r}{r}q^m=p^{-(r+1)}.
      \]
\end{enumerate}
Here \(q^0\defeq1\), also when \(q=0\).
Part~(2) with \(r=0\) is the geometric series
\(\sum_{m\geq0}q^m=1/p\).
\end{hblemma}

\section{The Poisson calculation}
\label{sec:p15-w02b-p15-001-poisson}

\begin{hbtheorem}{thm:p15-w02b-p15-001-poisson-moments}{proved=annex:p15-w02b-p15-001-poisson-moments}{Poisson factorial moments}
This is Durrett's Example~1.6.13 \cite[p.~34]{DurrettPTEAuthor2019}.
Let \(X\) have the Poisson distribution with parameter \(\lambda>0\).
For every integer \(k\geq1\), the descending product
\(X(X-1)\cdots(X-k+1)\) is nonnegative a.s. and
\[
 \mathbb E\bigl(X(X-1)\cdots(X-k+1)\bigr)=\lambda^k.
\]
Consequently \(\mathbb EX=\lambda\), \(\mathbb E[X(X-1)]=\lambda^2\),
\(\mathbb EX^2=\lambda^2+\lambda<\infty\), and
\[
 \operatorname{Var}(X)=\mathbb EX^2-(\mathbb EX)^2
   =\mathbb E[X(X-1)]+\mathbb EX-\lambda^2=\lambda.
\]
For the point mass at zero every factorial moment and the variance are
zero, which is the value \(\lambda^k\) at \(\lambda=0\).
\end{hbtheorem}

The orders one and two, and the variance, are also the moment assertions
of \cref{thm:p15-stirling-poisson-premises}; the theorem here adds every
higher order by the source's single index shift.

\section{The geometric calculation}
\label{sec:p15-w02b-p15-001-geometric}

\begin{hbtheorem}{thm:p15-w02b-p15-001-geometric-moments}{proved=annex:p15-w02b-p15-001-geometric-moments}{Geometric moments and first occurrence}
Parts~(1)--(3) are Durrett's Example~1.6.14
\cite[p.~35]{DurrettPTEAuthor2019}, with part~(2) extended to every order.
Let \(0<p<1\) and \(q=1-p\).
\begin{enumerate}
\item \(\sum_{k\geq1}pq^{k-1}=1\).
\item If \(N\) has the geometric distribution with success probability
      \(p\), then for every integer \(r\geq1\) the descending product
      \(N(N-1)\cdots(N-r+1)\) is nonnegative a.s. and
      \[
       \mathbb E\bigl(N(N-1)\cdots(N-r+1)\bigr)
        =\frac{r!\,q^{r-1}}{p^r}.
      \]
\item In particular
      \[
       \mathbb EN=\frac1p,\qquad
       \mathbb E[N(N-1)]=\frac{2(1-p)}{p^2},
      \]
      and
      \[
       \operatorname{Var}(N)=\frac{2(1-p)}{p^2}+\frac1p-\frac1{p^2}
        =\frac{1-p}{p^2}.
      \]
\item Let \(A_1,A_2,\ldots\) be independent events with
      \(\mathbb P(A_i)=p\) for every \(i\). Put
      \(N(\omega)\defeq\min\{i\geq1:\omega\in A_i\}\), and
      \(N(\omega)\defeq\infty\) when \(\omega\) lies in no \(A_i\).
      Then \(N\) is an extended real random variable,
      \(\mathbb P(N=\infty)=0\), and
      \(\mathbb P(N=k)=pq^{k-1}\) for every \(k\geq1\).
\end{enumerate}
\end{hbtheorem}

Part~(4) makes precise Durrett's reading of \(N\) as the count of
independent trials up to and including the first occurrence of an event
of probability \(p\)
\cite[p.~35]{DurrettPTEAuthor2019}. Changing \(N\) to the value one on
the null event \(\{N=\infty\}\) gives a real random variable with the
same law, to which parts~(2) and~(3) apply. At \(p=1\) the variable is
one a.s.; part~(2) then reads \(1\) for \(r=1\) and \(0\) for
\(r\geq2\), which is the displayed formula with the convention
\(0^0=1\). The orders one and two also appear, by a differentiation
argument, inside the proof of \cref{thm:p15-coupon-collector-law}.

\section{Procedure, worked values and limits}
\label{sec:p15-w02b-p15-001-procedure}

\begin{hbmethod}{def:p15-w02b-p15-001-method}{Computing a moment of a variable on the nonnegative integers}
\par\nobreak
\begin{enumerate}
\item Check that the variable takes values in \(\{0,1,2,\ldots\}\)
      with probability one. This is the side condition that makes each
      descending product nonnegative and makes the terms below its order
      vanish.
\item Write the factorial moment as a nonnegative series by part~(1) of
      \cref{thm:p15-w02b-p15-001-series-tools}. Discard the terms
      \(j<k\), replace the descending product by \(j!/(j-k)!\), and
      shift the index. A nonnegative series may be reindexed, regrouped
      and multiplied by a finite nonnegative constant without any
      convergence check (\cref{thm:p15-countable-atomic}).
\item Recover ordinary moments from the pointwise identities
      \(x^2=x(x-1)+x\) and \(x^3=x(x-1)(x-2)+3x(x-1)+x\), whose
      summands are nonnegative on the nonnegative integers, so that
      nonnegative additivity applies.
\item Subtract only finite quantities. The variance formula
      \[
       \operatorname{Var}(X)=\mathbb E[X(X-1)]+\mathbb EX-(\mathbb EX)^2
      \]
      is used only after \(\mathbb EX^2<\infty\) has been shown
      (\cref{thm:p15-expectation-bounds}).
\end{enumerate}
\end{hbmethod}

\begin{hbexample}{def:p15-w02b-p15-001-worked}{A third Poisson moment and a fair-coin waiting time}
Let \(X\) be Poisson with \(\lambda=2\). The theorem gives
\(\mathbb EX=2\), \(\mathbb E[X(X-1)]=4\) and
\(\mathbb E[X(X-1)(X-2)]=8\). Expanding,
\(x(x-1)(x-2)+3x(x-1)+x=x^3-3x^2+2x+3x^2-3x+x=x^3\). Each summand
is nonnegative on the nonnegative integers, so nonnegative additivity
gives
\[
 \mathbb EX^3=8+3\cdot4+2=22.
\]
Also \(\operatorname{Var}(X)=4+2-4=2\).

Let \(N\) be geometric with \(p=1/2\), the law of
\cref{def:p15-countable-waiting}. Then
\(\mathbb EN=2\), \(\mathbb E[N(N-1)]=2\cdot(1/2)/(1/4)=4\),
\(\mathbb EN^2=6\) and \(\operatorname{Var}(N)=6-4=2=q/p^2\).
As a direct check of the mean, the partial sums of
\(\sum_{k\geq1}k2^{-k}\) are
\(1/2,\ 1,\ 11/8,\ 13/8,\ 57/32\) for one to five terms; the
annex shows that the \(K\)-th partial sum is \(2-(K+2)2^{-K}\).
\end{hbexample}

\begin{hbwarning}{def:p15-w02b-p15-001-limits}{Where the calculation fails}
\par\nobreak
\begin{enumerate}
\item \emph{Values off the nonnegative integers.} If \(X=1/2\) a.s.,
      then \(X(X-1)=-1/4\) a.s. The descending product is negative and
      no term of the series vanishes for being below the order, so the
      index shift of the method has nothing to act on. The factorial
      moment route requires the side condition of step~(1).
\item \emph{An infinite mean.} Let \(M\) take the value \(k\geq1\) with
      probability \(1/(k(k+1))\). These masses sum to one, but
      \(\mathbb EM=\sum_{k\geq1}1/(k+1)=\infty\) and
      \(\mathbb E[M(M-1)]=\infty\). The expression
      \(\mathbb E[M(M-1)]+\mathbb EM-(\mathbb EM)^2\) would subtract an
      infinite quantity from an infinite quantity, and \(M\) has no
      variance since \(\mathbb EM^2=\infty\). Step~(4) of the method
      blocks this computation.
\item \emph{Success probability zero.} At \(p=0\) every geometric mass
      \(p(1-p)^{k-1}\) is zero, so the formula defines no probability
      law. In part~(4) of the geometric theorem with
      \(\mathbb P(A_i)=0\), the first-occurrence index equals \(\infty\)
      almost surely. The restriction \(0<p\) is therefore needed.
\end{enumerate}
These three assertions, and the partial sums of the worked example, are
proved at the end of the annex item for
\cref{thm:p15-w02b-p15-001-geometric-moments}.
\end{hbwarning}

\emph{Relationships.} The integral, the law formula and the variance
identities are \cref{thm:p15-expectation-law-integral,thm:p15-expectation-bounds}.
The Poisson normalization and the Poisson sums of independent
variables are in \cref{thm:p15-stirling-poisson-premises}; the
Poisson count enters the normal approximation in
\cref{thm:p15-clt-numerical-examples}. Geometric durations are the
building blocks of the coupon collector's time in
\cref{thm:p15-coupon-collector-law}. A variable whose signed mean
does not exist is in \cref{def:p15-computed-expectations}.

\emph{Sources.} Examples~1.6.13 and~1.6.14 of the author version dated
11 January 2019 \cite[pp.~34--35]{DurrettPTEAuthor2019}; the
appendix justification cited by the source is its Example~A.5.6
\cite[pp.~472--473]{DurrettPTEAuthor2019}.
